\documentclass[10pt, a4paper]{article}
\usepackage[utf8]{inputenc}
\usepackage[T1]{fontenc}
\usepackage{lmodern}
\usepackage{geometry}
\usepackage{amsmath}
\usepackage{amsfonts}
\usepackage{amssymb}
\usepackage{graphicx}
\usepackage{tikz}  % For AI generated diagrams
\usepackage{multicol} % For 2-column layout
\usepackage{fancyhdr}
\usepackage{enumitem}
\usepackage{float}
\usepackage{xcolor}  % For colored text

% TikZ Libraries (Standard set for Physics/Maths)
\usetikzlibrary{arrows.meta, calc, decorations.pathmorphing, patterns, shapes.geometric}

% Page Layout (JEE Standard)
\geometry{top=0.75in, bottom=1in, left=0.5in, right=0.5in}
\setlength{\columnsep}{0.6cm}

% Header and Footer
\pagestyle{fancy}
\fancyhf{}
\rhead{\textbf{Mentors Mantra}}
\lhead{\VAR{subject}}
\cfoot{
    \footnotesize
    \textcolor{gray}{\rule{0.4\textwidth}{0.4pt}} \\[2pt]
    \textcolor{gray}{For Free Trial Visit: \textbf{www.mentorsmantra.co.in}} \\[1pt]
    \textcolor{gray}{Contact: \textbf{9821040290} / \textbf{7982387231}} \\[1pt]
    \textcolor{gray}{\textbf{t.me/mentorsmantraofficial} -- Join us on Telegram for free guidance from IITians/NEET Experts}
}
\rfoot{\thepage}
\renewcommand{\footrulewidth}{0pt}

% Custom Option List (vertical line-by-line for long statements)
\newcommand{\fourchoices}[4]{%
    \vspace{2pt}%
    \par\noindent (a) #1%
    \par\noindent (b) #2%
    \par\noindent (c) #3%
    \par\noindent (d) #4%
    \vspace{4pt}%
}

\begin{document}

% --- TITLE SECTION ---
\begin{center}
    {\Huge \textbf{Mentors Mantra}} \\[0.2cm]
    {\Large \textbf{Target: \VAR{level}}} \\[0.05cm]
    {\normalsize \textbf{Difficulty: \VAR{difficulty}}} \\[0.1cm]
    \rule{\textwidth}{1pt}
    \textbf{Topic: \VAR{topic}}
\end{center}

\vspace{0.2cm}
\BLOCK{ if level == 'JEE Advanced' }
\textit{\textbf{Instructions:} This paper contains \VAR{total_questions} questions. Some MCQs are marked \textbf{(One or More Correct)} - select all correct options. Remaining MCQs have single correct answer. Section B contains Numerical Value Questions.}
\BLOCK{ else }
\textit{\textbf{Instructions:} This paper contains \VAR{total_questions} questions. Section A contains MCQs and Section B contains Numerical Value Questions.}
\BLOCK{ endif }
\vspace{0.3cm}

% --- QUESTIONS SECTION (Single Column) ---
\begin{enumerate}[label=\textbf{Q.\arabic*}, itemsep=10pt]

    % JINJA LOOP START
    \BLOCK{ for q in questions }

        \item \VAR{q.text}
        
        % Check if diagram exists
        \BLOCK{ if q.diagram_tikz }
        \begin{center}
        \begin{tikzpicture}[scale=0.6]
            \VAR{q.diagram_tikz}
        \end{tikzpicture}
        \end{center}
        \BLOCK{ endif }

        % Check if MCQ (single or multi), Numerical, or CBSE types
        \BLOCK{ if q.type == 'mcq_multi' }
            \textit{(One or More Correct)}
            \fourchoices{\VAR{q.options[0]}}{\VAR{q.options[1]}}{\VAR{q.options[2]}}{\VAR{q.options[3]}}
        \BLOCK{ elif q.type == 'mcq' }
            \fourchoices{\VAR{q.options[0]}}{\VAR{q.options[1]}}{\VAR{q.options[2]}}{\VAR{q.options[3]}}
        \BLOCK{ elif q.type == 'short_answer' }
            \hfill \textbf{[\VAR{q.marks|default(2)} Marks]}
            \vspace{0.3cm}
            \par\noindent\textcolor{gray}{\dotfill}
            \vspace{0.2cm}
            \par\noindent\textcolor{gray}{\dotfill}
            \BLOCK{ if q.marks|default(2) >= 3 }
            \vspace{0.2cm}
            \par\noindent\textcolor{gray}{\dotfill}
            \BLOCK{ endif }
        \BLOCK{ elif q.type == 'long_answer' }
            \hfill \textbf{[5 Marks]}
            \vspace{0.3cm}
            \par\noindent\textcolor{gray}{\dotfill}
            \vspace{0.2cm}
            \par\noindent\textcolor{gray}{\dotfill}
            \vspace{0.2cm}
            \par\noindent\textcolor{gray}{\dotfill}
            \vspace{0.2cm}
            \par\noindent\textcolor{gray}{\dotfill}
            \vspace{0.2cm}
            \par\noindent\textcolor{gray}{\dotfill}
        \BLOCK{ elif q.type == 'case_based' }
            \hfill \textbf{[4 Marks]}
            \vspace{0.2cm}
            \par\textit{\VAR{q.passage|default('')}}
            \vspace{0.3cm}
            \BLOCK{ if q.sub_questions }
            \BLOCK{ for sq in q.sub_questions }
            \par\noindent (\VAR{loop.index|string|replace('1','i')|replace('2','ii')|replace('3','iii')|replace('4','iv')}) \VAR{sq.text}
            \BLOCK{ if sq.options }
            \fourchoices{\VAR{sq.options[0]}}{\VAR{sq.options[1]}}{\VAR{sq.options[2]}}{\VAR{sq.options[3]}}
            \BLOCK{ endif }
            \BLOCK{ endfor }
            \BLOCK{ endif }
        \BLOCK{ else }
            \BLOCK{ if q.marks }
            \hfill \textbf{[\VAR{q.marks} Marks]}
            \BLOCK{ endif }
            \vspace{0.2cm} \\ 
            \textbf{Answer: \_\_\_\_\_\_\_\_\_}
        \BLOCK{ endif }

    \BLOCK{ endfor }
    % JINJA LOOP END

\end{enumerate}

\newpage

% --- ANSWER KEY SECTION ---
\begin{center}
    {\Large \textbf{ANSWER KEY}}
\end{center}

\begin{table}[h]
\centering
\renewcommand{\arraystretch}{1.5}
\begin{tabular}{|c|c||c|c||c|c|}
\hline
\textbf{Q.No} & \textbf{Ans} & \textbf{Q.No} & \textbf{Ans} & \textbf{Q.No} & \textbf{Ans} \\
\hline
\BLOCK{ for q in questions }
    \textbf{\VAR{loop.index}} & \VAR{q.answer} 
    \BLOCK{ if loop.index is divisibleby 3 } \\ \hline \BLOCK{ else } & \BLOCK{ endif }
\BLOCK{ endfor }
\\ \hline
\end{tabular}
\end{table}


% --- SOLUTIONS SECTION (Step-by-Step Format) ---
\BLOCK{ if include_solutions }
\newpage

\begin{center}
    {\Large \textbf{Detailed Solutions}}\\[0.1cm]
    {\normalsize \VAR{subject} -- \VAR{topic}}
\end{center}
\vspace{0.5cm}

\BLOCK{ for q in questions }
\noindent\rule{\textwidth}{0.4pt}
\vspace{0.3cm}

\noindent{\large \textbf{Question \VAR{loop.index}.}} \hfill \fbox{\textbf{Answer: \VAR{q.answer}}}

\vspace{0.3cm}
\noindent\textit{Q: \VAR{q.text}}

\vspace{0.4cm}
\noindent\textbf{Solution:}
\vspace{0.2cm}

\BLOCK{ if q.solution }
\VAR{q.solution}
\BLOCK{ else }
The correct option is \textbf{\VAR{q.answer}}.
\BLOCK{ endif }

\vspace{0.5cm}
\BLOCK{ endfor }

\BLOCK{ endif }

\end{document}
